\section{Experimental Setup}
\label{sec:experimental_setup}

\subsection{Data and common evaluation protocol}
Table~\ref{tab:datasets} lists eight multivariate series: ETTh1, ETTh2, ETTm1, and ETTm2 from the Electricity Transformer Temperature data~\cite{ettdata}; Appliances Energy Prediction~\cite{appliancesdata}; and the Rat, Fox, and Panther sites from Building Data Genome 2 (BDG2)~\cite{bdg2data}. Each ETT series has seven channels. For Appliances, we omit the timestamp and the random variables \texttt{rv1} and \texttt{rv2}, leaving 26 numerical channels. Each BDG2 site contains 15 selected electricity meters. Forecasts and losses include every retained channel.

The BDG2 meter IDs and channel order are fixed in the released inputs. Meters were chosen by their missing-value fractions over the supplied series. Fox and Panther were forward filled and then backward filled from the pinned cleaned source. The archived Rat input differs from that procedure at 78 cells; explicit snapshot adjustments reproduce its evaluated values.

For every series, the lookback is $L=96$ samples, the horizon is $H=24$ samples, and consecutive forecast origins are 24 samples apart. These sample counts represent different durations: ETTh and BDG2 are hourly, ETTm is sampled every 15 minutes, and Appliances every 10 minutes. Let $T$ be the series length, $n=\lfloor T/2\rfloor$, and $a=\lfloor0.8n\rfloor$. The legacy base is fitted using $[0,a)$, with $[a,n)$ reserved for model and comparator selection. A refit base uses $[0,n)$ after its settings have been selected. Evaluation origins are $n,n+H,\ldots$; only complete target blocks within $[n,T)$ are scored. Per-channel means and standard deviations are computed on $[0,n)$; we subtract the mean and divide by the standard deviation plus $10^{-8}$. All reported errors use this normalized scale.

\begin{table}[!t]
\caption{Prepared evaluation series. Eval. is the zero-based index of the first evaluation target; Blocks counts complete 24-sample targets. All eight series enter the same fixed-base comparison.}
\label{tab:datasets}
\centering
% Dataset-specific quantities. Common L, H, and stride are defined in Section IV.
\begingroup\fontsize{9}{11}\selectfont
\setlength{\tabcolsep}{2.3pt}
\begin{tabular}{lrrrr}
\toprule
Series & Samples & $C$ & Eval. & Blocks \\
\midrule
ETTh1 & 17420 & 7 & 8710 & 362 \\
ETTh2 & 17420 & 7 & 8710 & 362 \\
ETTm1 & 69680 & 7 & 34840 & 1451 \\
ETTm2 & 69680 & 7 & 34840 & 1451 \\
Appliances & 19735 & 26 & 9867 & 411 \\
BDG2 (Rat) & 17544 & 15 & 8772 & 365 \\
BDG2 (Fox) & 17544 & 15 & 8772 & 365 \\
BDG2 (Panther) & 17544 & 15 & 8772 & 365 \\
\bottomrule
\end{tabular}
\endgroup

\end{table}

The fixed-base comparison uses DLinear~\cite{dlinear2022} and a compact, channel-independent PatchTST~\cite{patchtst2023}. DLinear decomposes the input with a replicate-padded moving average of width 25 and uses separate trend and remainder maps shared across channels. PatchTST uses length-16 patches at stride 8, a 64-dimensional embedding, four attention heads, two encoder layers, a feed-forward width of 128, zero dropout, a learned position embedding, and a linear forecast head. It also uses a learned affine input transformation per channel.

For each series, architecture, and seed in $\{0,1,2\}$, we train the legacy base with Adam at $10^{-3}$, batch size 32, and MSE loss on randomly sampled training windows. The update count is chosen by validation MSE from $\{200,500,1000,2000,4000,8000,20000\}$. The refit base is reinitialized with the corresponding seed and trained on $[0,n)$ for that selected count. Both models are frozen during evaluation. The eight series, two architectures, three seeds, and two training variants yield $8\times2\times3\times2=96$ conditions. Within each condition, all methods receive the same base forecasts, observed targets, and forecast origins. Each method issues a complete forecast before the target arrives and updates only after the full target block is observed. The first evaluation block is included in every reported loss.

\subsection{Comparators and learned comparison design}
Table~\ref{tab:prior_methods} distinguishes the auxiliary inputs and retained states of related methods. The main fixed-base comparison contains seven arms: the uncorrected base (Static), the input/output $\delta$-Adapter~\cite{deltaadapter2026}, COSA~\cite{cosa2026}, FAC~\cite{fac2026}, OMPB~\cite{ompb2026}, full ELF~\cite{elf2025}, and EPOC with $K=4$. All six correction arms use the same bounded, global, 32-completed-block blending rule of Section~\ref{sec:method}. Online states start empty at the evaluation boundary; weights trained before evaluation are retained. The base remains fixed, including when FAC calibrates its input.

The $\delta$-Adapter has bounded additive input and output multilayer perceptrons of width 512. The additive bound is 0.01 on ETT and 0.1 on the other series. For each fixed base, the adapter receives ten source-training epochs with Adam on at most 1024 sampled pre-validation windows. It takes one clipped Adam step when each evaluation target completes. Its online rate is selected from $\{10^{-6},10^{-4},10^{-3}\}$ using chronological online MSE on legacy validation origins starting at $a+96$. COSA uses the official linear adapter and adaptive rate function from commit \texttt{527c0feb}, with channel-specific maps and gates and ten entries of completed-block context. FAC is a reconstruction of the published input/output frequency-calibration equations using channel-specific complex masks and biases, inverse real FFT, and a gated correction. COSA and FAC take three gradient steps per completed block. Their initial rates are selected on the same legacy validation origins from $\{10^{-4},10^{-3},5\times10^{-3},10^{-2}\}$. Each selection uses that series, architecture, and seed; the chosen rate is carried to the corresponding refit condition, and adapter state is reset before evaluation.

OMPB uses a gated per-channel Gaussian residual head with source-risk, divergence, and supervised-feedback terms. Its transferred settings are $K=5$, five optimizer steps per completed block, Adam rate 0.005, a 256-pair source buffer, and a 64-forecast target buffer. The head is pretrained for 20 epochs using at most 1024 source windows per fixed base; no validation search changes these settings. Full ELF uses cumulative Fourier-coordinate ridge regression with one newly completed training window per block, no pre-evaluation regression pairs, and ridge coefficient 20. Its input/output frequency counts are $(d,q)=(87,10)$. The input is centered by its channel mean before orthonormal Fourier transformation; the mean is restored after reconstruction. In unscaled sufficient statistics, the channel-$c$ ridge term is $20\sigma_c^2 I$. Before fitting, the auxiliary forecast repeats the most recent daily cycle. We retain independent real Fourier coordinates and packed symmetric statistics for storage accounting. ELF's original dynamic weighting is replaced by the shared global controller in the main comparison.

A separate experiment compares EPOC with the two-projection TEFL-style residual adapter~\cite{tefl2026}. It covers six series (the four ETT series, Appliances, and Rat), both architectures, and three seeds. The adapter has shared weights across channels, a ReLU hidden layer, no projection biases, and widths 2, 4, or 64; its second projection is initialized to zero. Base warmup lasts three chronological epochs under either MAE alone or MAE plus spectral flatness (SF) with unit weight. SF is the geometric-to-arithmetic mean ratio of full Fourier power, stabilized by $10^{-8}$, after consecutive residual blocks in a batch are concatenated in time and averaged over channels. The second stage uses 12 epochs of MAE optimization, either jointly for base and adapter or for the base alone. Both stages use AdamW at $10^{-3}$, weight decay 0.01, and batch size 32. Checkpoints are selected by validation MAE over epochs 0--12, with earlier epochs winning ties. Paired runs share the warmup state and shuffle seed.

At evaluation, a learned adapter receives the previous completed base residual and has fixed weights. We score its raw and globally blended forecasts. For each warmup and width, the direct EPOC-versus-adapter comparison applies both corrections to the \emph{same jointly trained base} with the same blending rule. A separately trained base-only run tests the benefit of the joint adapter procedure. We aggregate this group separately from the 96 fixed-base conditions.

The endpoint-location control replaces the most recent residual in EPOC's regression with the first residual, the twelfth residual, or the mean residual of the preceding block. Each scalar is multiplied by $\sqrt H$, enters all $K$ regressions of its channel, and has the same retained size as the endpoint. We run this control on all 96 fixed-base conditions and, separately, on the width-64 jointly trained bases under each warmup. For the fixed-base group, input ablations additionally remove the endpoint, the residual DCT coefficients, or the current forecast coefficients; compare endpoint-only regression and endpoint persistence; and supply an endpoint projected from the retained residual DCT coefficients. All fitted variants retain the same $K$-component output subspace. A 13-setting one-factor sweep changes $K$ to $1,2,3,6,8$, ridge strength to $0.1$ or $10$, forgetting half-life to $32,512$, or infinity completed blocks, and blending window to $8$ or $128$ blocks, alongside the reference setting. The input and setting studies each use all 96 fixed-base conditions.

\subsection{Metrics and aggregation}
For $N$ completed evaluation blocks, write $E_{t,h,c}=\widehat y_{t,h,c}-y_{t,h,c}$. We report normalized-scale errors
\begin{equation}
 \begin{aligned}
 \operatorname{MSE}&=\frac{1}{NHC}\sum_{t,h,c}E_{t,h,c}^{2},\\
 \operatorname{MAE}&=\frac{1}{NHC}\sum_{t,h,c}|E_{t,h,c}|.
 \end{aligned}
 \label{eq:eval_metrics}
\end{equation}
For metric $m\in\{\mathrm{MSE},\mathrm{MAE}\}$ and condition $i$, the reduction from Static is $100(1-m_{i,\mathrm{method}}/m_{i,\mathrm{Static}})$. A paired reduction from comparator $b$ to method $p$ instead uses $100(1-m_{i,p}/m_{i,b})$; positive values favor $p$. These are different denominators, and MSE and MAE are calculated separately.

Main overview values average the 96 condition-wise reductions and take the median of retained numerical-array bytes over those conditions. Each point in Fig.~\ref{fig:accuracy_storage} summarizes the three seeds and two training variants for one series and one architecture. In the jointly trained experiment, we first average each method's losses over three seeds within each of 12 series--architecture pairs, then compute the paired reduction and average those 12 percentages. Endpoint-location results keep the fixed-base and jointly trained groups separate. The order of aggregation matters: a mean of paired percentages generally differs from a percentage calculated from pooled losses.
